% Options for packages loaded elsewhere
\PassOptionsToPackage{unicode}{hyperref}
\PassOptionsToPackage{hyphens}{url}
\PassOptionsToPackage{dvipsnames,svgnames,x11names}{xcolor}
%
\documentclass[
  letterpaper,
  DIV=11,
  numbers=noendperiod]{scrartcl}

\usepackage{amsmath,amssymb}
\usepackage{iftex}
\ifPDFTeX
  \usepackage[T1]{fontenc}
  \usepackage[utf8]{inputenc}
  \usepackage{textcomp} % provide euro and other symbols
\else % if luatex or xetex
  \usepackage{unicode-math}
  \defaultfontfeatures{Scale=MatchLowercase}
  \defaultfontfeatures[\rmfamily]{Ligatures=TeX,Scale=1}
\fi
\usepackage{lmodern}
\ifPDFTeX\else  
    % xetex/luatex font selection
\fi
% Use upquote if available, for straight quotes in verbatim environments
\IfFileExists{upquote.sty}{\usepackage{upquote}}{}
\IfFileExists{microtype.sty}{% use microtype if available
  \usepackage[]{microtype}
  \UseMicrotypeSet[protrusion]{basicmath} % disable protrusion for tt fonts
}{}
\makeatletter
\@ifundefined{KOMAClassName}{% if non-KOMA class
  \IfFileExists{parskip.sty}{%
    \usepackage{parskip}
  }{% else
    \setlength{\parindent}{0pt}
    \setlength{\parskip}{6pt plus 2pt minus 1pt}}
}{% if KOMA class
  \KOMAoptions{parskip=half}}
\makeatother
\usepackage{xcolor}
\setlength{\emergencystretch}{3em} % prevent overfull lines
\setcounter{secnumdepth}{-\maxdimen} % remove section numbering
% Make \paragraph and \subparagraph free-standing
\ifx\paragraph\undefined\else
  \let\oldparagraph\paragraph
  \renewcommand{\paragraph}[1]{\oldparagraph{#1}\mbox{}}
\fi
\ifx\subparagraph\undefined\else
  \let\oldsubparagraph\subparagraph
  \renewcommand{\subparagraph}[1]{\oldsubparagraph{#1}\mbox{}}
\fi


\providecommand{\tightlist}{%
  \setlength{\itemsep}{0pt}\setlength{\parskip}{0pt}}\usepackage{longtable,booktabs,array}
\usepackage{calc} % for calculating minipage widths
% Correct order of tables after \paragraph or \subparagraph
\usepackage{etoolbox}
\makeatletter
\patchcmd\longtable{\par}{\if@noskipsec\mbox{}\fi\par}{}{}
\makeatother
% Allow footnotes in longtable head/foot
\IfFileExists{footnotehyper.sty}{\usepackage{footnotehyper}}{\usepackage{footnote}}
\makesavenoteenv{longtable}
\usepackage{graphicx}
\makeatletter
\def\maxwidth{\ifdim\Gin@nat@width>\linewidth\linewidth\else\Gin@nat@width\fi}
\def\maxheight{\ifdim\Gin@nat@height>\textheight\textheight\else\Gin@nat@height\fi}
\makeatother
% Scale images if necessary, so that they will not overflow the page
% margins by default, and it is still possible to overwrite the defaults
% using explicit options in \includegraphics[width, height, ...]{}
\setkeys{Gin}{width=\maxwidth,height=\maxheight,keepaspectratio}
% Set default figure placement to htbp
\makeatletter
\def\fps@figure{htbp}
\makeatother

\KOMAoption{captions}{tableheading}
\makeatletter
\@ifpackageloaded{caption}{}{\usepackage{caption}}
\AtBeginDocument{%
\ifdefined\contentsname
  \renewcommand*\contentsname{Table of contents}
\else
  \newcommand\contentsname{Table of contents}
\fi
\ifdefined\listfigurename
  \renewcommand*\listfigurename{List of Figures}
\else
  \newcommand\listfigurename{List of Figures}
\fi
\ifdefined\listtablename
  \renewcommand*\listtablename{List of Tables}
\else
  \newcommand\listtablename{List of Tables}
\fi
\ifdefined\figurename
  \renewcommand*\figurename{Figure}
\else
  \newcommand\figurename{Figure}
\fi
\ifdefined\tablename
  \renewcommand*\tablename{Table}
\else
  \newcommand\tablename{Table}
\fi
}
\@ifpackageloaded{float}{}{\usepackage{float}}
\floatstyle{ruled}
\@ifundefined{c@chapter}{\newfloat{codelisting}{h}{lop}}{\newfloat{codelisting}{h}{lop}[chapter]}
\floatname{codelisting}{Listing}
\newcommand*\listoflistings{\listof{codelisting}{List of Listings}}
\makeatother
\makeatletter
\makeatother
\makeatletter
\@ifpackageloaded{caption}{}{\usepackage{caption}}
\@ifpackageloaded{subcaption}{}{\usepackage{subcaption}}
\makeatother
\ifLuaTeX
  \usepackage{selnolig}  % disable illegal ligatures
\fi
\usepackage{bookmark}

\IfFileExists{xurl.sty}{\usepackage{xurl}}{} % add URL line breaks if available
\urlstyle{same} % disable monospaced font for URLs
\hypersetup{
  colorlinks=true,
  linkcolor={blue},
  filecolor={Maroon},
  citecolor={Blue},
  urlcolor={Blue},
  pdfcreator={LaTeX via pandoc}}

\author{}
\date{}

\begin{document}

\subsection{Introduccion}\label{introduccion}

\subsubsection{Motivacion}\label{motivacion}

Queremos Analizar los datos de la actividad de vigilancia entomológica
de forma oportuna y automática utilizando~herramientas de Ciencia de
Datos utilizando la base de datos de la plataforma de ``Vigilancia
Entomológica y Control Integral del Vector'' del INSTITUTO NACIONAL DE
SALUD PÚBLICA~en colaboración con Secretaría de Salud y el Centro
Nacional de Prevencion y Control de enfermedades (CENAPRECE). Con el fin
de calcular y analizar los índices de riesgo entomológico, y con esto
gestionar información para la toma de decisiones en la prevención de
enfermedades transmitidas por vector.

El análisis de riesgo entomológico se hace en hojas de cálculo El manejo
de la información es ineficiente: La estructura de la base de datos es
compleja Adaptar consultas a pequeñas modificaciones es impráctico La
visualización georreferenciada de esta plataforma es pobre e inestable

Usar R para visualizar la información en tiempo y espacio Automatizar el
manejo de la información Generar visualización georreferenciada con los
índices de estegomía

\subsubsection{Descripción del
problema}\label{descripciuxf3n-del-problema}

El Programa estatal de vigilancia entomología y control integral de
enfermedades trasmitidas por vector, realiza diversas actividades para
la prevención de enfermedades, las cuales se guían con actividades de
vigilancia que proporcionan índices de riesgo entomológico, por ello se
hará una API para calcular los índices de riesgo entomológico de forma
automática y oportuna para la toma de decisione.

Todas las actividades se capturan diariamente en una plataforma, donde
es posible descargar los datos para su reporte y análisis. En el caso de
Vigilancia por Ovitrampa(VO), la plataforma calcula los índices de
riesgo pero para la actividad de Estudio Entomológico (EE) los calcula
por manzana trabajada, estos índices se requiere por localidad de riesgo
de manera semanal.

Se requiere poder calcular índices de riesgo del EE por tipo de estudio,
localidad de riesgo y por semana epidemiológica de manera automatizada.

La actividad de EE evalúa una muestra antes(encuesta) y
después(verificación) de las actividades de control integral para
determinar el riesgo entomológico\\
que hay en un área grande a trabajar (Sectores o Localidad) y si
requiere reforzar acciones de control y prevención.

Tener la información oportuna de los índices de riego entomológicos
permiten tomar decisiones informades de ¿En Dónde? y ¿Cuándo? focalizar
intervenciones.

\subsubsection{Gestion de datos}\label{gestion-de-datos}

Para validar y poner a prueba la solución propuesta, se plantea la
utilización de una base de datos ficticia que sea representativa con las
bases de datos reales descargadas de la plataforma. Esto permitirá
evaluar el funcionamiento del paquete de R en un entorno controlado. Una
vez realizadas las pruebas iniciales y asegurada la confiabilidad de los
resultados obtenidos,se establecerá un convenio de confidencialidad con
el fin de acceder a los datos reales y evaluar el desempeño del paquete
de R en un escenario más cercano a la realidad. Esta estrategia
garantizará la confidencialidad de los datos utilizados, a la vez que
brindará la oportunidad de comprobar la efectividad de la solución en
situaciones reales y aplicables al contexto de vigilancia entomológica y
control de enfermedades transmitidas por vectores.



\end{document}
